\honest{How to Convince Me That 2 + 2 = 3}{Eliezer Yudkowsky}{2007}

In ``What Is Evidence?'' I wrote that a belief is only worthwhile if you could, in principle, be persuaded to believe otherwise. Cihan Baran replied: ``I can not conceive of a situation that would make 2 + 2 = 4 false. Perhaps for that reason, my belief in 2 + 2 = 4 is unconditional.'' \nb{The same comment also objected that mathematical evidence is not an event, since a theorem ``doesn't cause anything.'' The post does not quote that objection.}

I cannot conceive of such a situation either. (Redefining the symbols is not a situation; then you are no longer talking about 2, 4 or +.) But what could make a fact false and what could change my belief in it are different questions.

Here is what would convince me that 2 + 2 = 3. One morning I set two earplugs beside two others and find three, with none having ``appeared or disappeared.'' When I picture xx and xx, I get xxx. It checks with subtraction: xx from xxx leaves xx, while xx from xxxx would leave xxx. My memory would seem absurd against that. A calculator, Google and my copy of \textsc{1984} all agree. \nb{In the novel Winston writes that freedom is ``the freedom to say that two plus two make four.'' The change is part of the joke.} I would explain my old belief as a memory fault, perhaps from a sneeze, or as someone tampering with me, by hypnosis or because I am a simulation. Even then, I would think it likelier they had tampered with my memory of arithmetic than that 2 + 2 really equals 4. Either way I would notice that I was very confused. \nb{Both explanations are about the author's own mind. On the view that arithmetic is known a priori, such evidence can undercut belief even in a proof. So the scenario shows that the belief could change, not that it rests on experience.}

This is ``exactly the same kind of evidence'' that now convinces me that 2 + 2 = 4: ``physical observation, mental visualization, and social agreement.'' \nb{That is a disputed position, close to John Stuart Mill's, and the post states it as plain fact. Picturing units one after another is Kant's example of how arithmetic is known a priori. Even Quine, who held that arithmetic could in principle be revised, held that no particular experience confirms 2 + 2 = 4.}

Once I did not know that 2 + 2 = 4, and I did not learn it at random. \nb{How a belief was acquired does not settle what justifies it. Frege said that the difference between a priori and empirical knowledge lies in ``the justification for making the judgement.''} My brain came to store an answer that matches what earplugs do, and that match needs explaining. A belief either got there by a process entangled with reality, or it is right only by coincidence. For any belief complex enough to need more than about 10 bits to describe, coincidence is out of the question.

So unconditional facts are not unconditional beliefs, and my belief is not blind faith. Entangled evidence could convince me that a fact is unconditional. I do not say how. \nb{Kant, holding arithmetic necessary, denied that it comes from experience. Mill, holding that it comes from experience, gave up its necessity. The post keeps both.}

In a footnote I ask Christians what would convince them of Islam. ``Presumably'' the answer is being born to a Muslim mother and raised by Muslim parents. If there is more to it, what would convince them, and would their kind of reasoning have freed them from their religion had they been raised Muslim? \nb{For the author, ``what would convince me'' meant evidence. For the Christian reader, the post supposes, it means upbringing.}
